\begin{papercard}{brackin2026}
\cf{Stages} \stagetag{C} (a definition of contradiction between a commitment and a practice in the same policy, close to Axis C's ``practice commitment'' and ``potential mismatch''), \stagetag{L} (multi-stage LLM pipeline with a three-model judge panel and a stability re-run). Single-author arXiv preprint of September 2026. Added from the gap search of 23 October 2026 for its concept and as a cautionary design.

\cf{Concept} A commitment such as ``we do not sell your personal information'' stands beside sharing practices described elsewhere \pc{brackin2026}{1}{A company’s privacy policy may promise “we do not sell your personal information” in one section while describing extensive third-party data sharing in another}. The contradiction is pragmatic, not logical: the broad reassurance is undermined by a specific practice \pc{brackin2026}{3}{do not target pragmatic contradictions where a broad reassurance is undermined by a specific practice}. In contextual-integrity terms, a stated norm is undermined by a documented flow \pc{brackin2026}{5}{privacy washing can be read in these terms as a stated norm undermined by a documented flow within the same document}. Commitments phrased with ``may'' are excluded because they imply no limit.

\cf{Pipeline} Statement extraction, compatibility filtering and natural-language-inference screening, a multi-model judge and thematic analysis \pc{brackin2026}{1}{This paper operationalizes this phenomenon, privacy washing, through a four-stage automated pipeline}; a contradiction counts when a majority of a three-model panel confirms it \pc{brackin2026}{1}{Contradictions are confirmed by majority vote of a three-model large language model (LLM) panel}.

\cf{Results} At least one confirmed contradiction in 12.2\,\% of 123 companies (2026 corpus) and 36.5\,\% of 115 (OPP-115, 2015) \pc{brackin2026}{1}{At least one panel-confirmed contradiction appears in 12.2\% of OPPT companies (15/123}. The most frequent pattern concerns third-party sharing, including ``we do not sell or share'' statements added for the CCPA beside older sharing descriptions \pc{brackin2026}{6}{The SALE_SHARING pattern is driven by CCPA-mandated disclosures}. A re-run seven months later with different models reproduced the 2026 prevalence (13.0\,\%), but the dominance of third-party sharing depended on the panel.

\cf{Limits (stated by the author)} No human validation: precision is unknown and prevalence is a lower bound \pc{brackin2026}{1}{panel verdicts have not been validated against human expert judgment, so precision is unknown, and the pipeline’s aggressive filtering makes the prevalence figures lower bounds}. The two corpora were run with different filters, so their difference is not interpretable. \textbf{Card:} website policies in English, US law; a three-model majority is a consistency check among models, not reliability, the same point Gilardi et al.\ \pp{gilardi2023} illustrate for self-agreement.

\cf{Use for this thesis} (1) Axis C vocabulary: a no-sale or ``we never share'' pledge is a \emph{practice commitment}; a disclosed flow that falls inside its scope makes a candidate \emph{potential mismatch}, which the methodology routes to human review. The hedging exclusion (``may'') matches the boundary between commitment and qualified wording. (2) Within-document comparison is a cheaper first step than cross-surface comparison (Ferrari et al.\ \pp{ferrari2026}). (3) A model-panel verdict without human validation is the design the thesis avoids; cite it as such when justifying manual review of all mismatches.
\end{papercard}
